% Restrictions on a compression method, not absolute optimality claims.
\section{What generic projection can and cannot compress}
\label{p3sys:projection-limits}

The finite-fiber method eliminates algebraic choices over a retained
base.  It cannot automatically replace freely varying base coordinates
by a single primitive element.  The distinction is geometric, and
it must not be confused with a lower bound for every possible
Diophantine definition.

\begin{proposition}\label{p3sys:finite-projection-dimension}
Let $X$ be an integral $K$-variety of dimension $d$.  If a rational
map from a nonempty open subset of $X$ to $\mathbb A^b_K$ has
finite geometric fibers, then $b\ge d$.
In particular, rational changes of coordinates followed by finite-fiber
projection cannot reduce the dimension of the retained base below $d$.
\end{proposition}
\begin{proof}
Let $Y$ be the closure of the image.  A generically finite map makes
the function-field extension $K(X)/K(Y)$ finite.  Their transcendence
degrees over $K$ are therefore equal.  Thus
$d=\dim Y\le b$.  Restricting to a nonempty open or making a
birational change does not change the function field or its
transcendence degree.
\end{proof}

\begin{corollary}\label{p3sys:monic-base-barrier}
Suppose a source is described by a positive-rank finite free algebra
over $K[X_1,\ldots,X_b]$, followed only by nonzero conditions.
Every nonempty irreducible component of its geometric open locus
has dimension $b$.  A compression which preserves such a component
up to birational equivalence and finite-fiber projection must retain
at least $b$ base coordinates.
\end{corollary}
\begin{proof}
A finite algebra is integral over its base, and every minimal
prime of a finite flat algebra over a domain lies over the zero
prime.  To see the latter assertion, use going down for a flat
extension: a nonzero contracted prime would admit a strictly
smaller prime below the proposed minimal prime.  The induced
finite dominant map from each irreducible component to the
base preserves dimension.  A nonempty open in that component
still has dimension $b$.  Apply the preceding proposition.
\end{proof}

The unstratified 43-coordinate source before the direct Lucas-residual
refinement has 42 retained base coordinates and a triangular monic
finite algebra, followed only by guards.  The new direct-residual
source has 40 retained coordinates and the same property.  Thus a
generic projection does not, by itself, turn either source into a
finite-fiber presentation over nine retained coordinates.  Such an
operation would have positive-dimensional geometric fibers.  A proof
that those fibers retain the required rational points would be new
arithmetic information, not a consequence of the primitive-element
theorem.  The subsequent decomposition into zero and nonzero quadratic
root cases can introduce closed conditions and smaller components;
the dimension assertion just made concerns the original unstratified
source, not each of those later cases separately.

\begin{proposition}\label{p3sys:no-generic-rational-section}
Let a source variety with regular input $T$ project to
$\mathbb A^b_K$ with coordinates $(T,Z_2,\ldots,Z_b)$.
If that projection has a rational section on a nonempty open
subset of $\mathbb A^b_K$, then the source's input projection
contains all but finitely many elements of $K$.
Consequently a source whose rational input projection is
$\mathcal O_K$ has no such section.
\end{proposition}
\begin{proof}
Shrink the domain of the section to a nonempty principal open
$G(T,\mathbf Z)\ne0$.  For all but finitely many $t\in K$,
the specialization $G(t,\mathbf Z)$ is a nonzero polynomial:
at least one of its nonzero coefficient polynomials in $T$
has only finitely many roots.  Since $K$ is infinite, one can
choose $\mathbf Z\in K^{b-1}$ with $G(t,\mathbf Z)\ne0$.
The section then supplies a rational source point at input $t$.
The ring $\mathcal O_K$ has infinite complement in $K$, so it
cannot contain such a cofinite projection.
\end{proof}

The role of the base in this statement matters.  A rational section
over an already constrained subvariety is compatible with an
arithmetic image.  For example, in the presentation
\[
                         Y^2=T
\]
of rational squares one may add arbitrarily many completely free
coordinates.  They can later be deleted with no change of rational
projection, although the deleting map has positive-dimensional fibers.
Its section is over the square-cover variety, not over the entire
input affine line.  Thus dimension alone is not a lower bound on the
number of coordinates in every equivalent presentation.

There is a second reason for caution.  The geometric dimension of
one equation need not be the dimension of the closure of its rational
solutions.  The equation
\[
                         U^2+V^2=0
\]
defines a one-dimensional hypersurface over an algebraic closure,
but its rational solution set is just $(0,0)$.  Therefore a large
geometric source can in principle have its rational points confined
to a much smaller algebraic set.  The dimension restriction must
be applied to a component actually preserved by the proposed
compression, not silently to all equivalent rational presentations.

\begin{corollary}\label{p3sys:arithmetic-dimension-floor}
For every polynomial-system definition of $\mathcal O_K$, the
Zariski closure of its rational witness points has a horizontal
component of dimension at least two.  On that component the input
map has positive-dimensional generic geometric fiber.
\end{corollary}
\begin{proof}
Use Lemma~\ref{p3sys:horizontal-closure}.  Its nonconstant input
function is dominant onto the affine line, so the generic fiber
has dimension $\dim X-1\ge1$.
\end{proof}

This last restriction is intrinsic to the rational-point locus,
unlike the preceding bound for a particular finite-flat presentation.
It rules out compressing all witnesses to finitely many algebraic
choices over the input alone.  It does not rule out a two-dimensional
arithmetic witness variety, nor a definition with three to ten total
coordinates.  No such new construction is furnished by these
projection arguments.
